\Askhsh{28476_SOLUTION}{1}
\begin{ylh}{1}
    \item [Ύλη από εικόνα]
\end{ylh}
\begin{alist}
    \item [α)]
    \begin{rlist}
        \item Είναι:
        \[ \lim_{x \to 1} \frac{\ln x}{x-1} = \lim_{x \to 1} \frac{(\ln x)'}{(x-1)'} = \lim_{x \to 1} \frac{\frac{1}{x}}{1} = \lim_{x \to 1} \frac{1}{x} = 1 \]
        \item Είναι $f(x) = \frac{\ln x}{x-1} \cdot \frac{f(x)(x-1)}{\ln x}$. Άρα
        \[ \lim_{x \to 1} f(x) = \lim_{x \to 1} \left(\frac{f(x)(x-1)}{\ln x}\right) \cdot \lim_{x \to 1} \left(\frac{\ln x}{x-1}\right) \]
        Επομένως έχουμε:
        \[ \lim_{x \to 1} f(x) = 0 \cdot 1 = 0 \]
        Η συνάρτηση $f$ είναι συνεχής στο $\mathbb{R}$ ως παραγωγίσιμη, επομένως είναι
        \[ \lim_{x \to 1} f(x) = f(1) \Rightarrow f(1) = 0. \]
        $2^{\text{ος}}$ τρόπος:

        Η συνάρτηση $f$ είναι συνεχής στο $\mathbb{R}$ ως παραγωγίσιμη, επομένως είναι
        \[ \lim_{x \to 1} f(x) = f(1). \]
        Αρκεί λοιπόν να αποδείξουμε ότι
        \[ \lim_{x \to 1} f(x) = 0. \]
        Θεωρούμε συνάρτηση $g$ με $g(x) = \frac{f(x)(x-1)}{\ln x}$, οπότε είναι
        \[ f(x) = \frac{g(x) \ln x}{x-1} \text{, } x \in (0,1) \cup (1,+\infty). \]
        Επειδή είναι
        \[ \lim_{x \to 1} g(x) = 0 \text{ και } \lim_{x \to 1} \frac{\ln x}{x-1} = \lim_{x \to 1} \frac{(\ln x)'}{(x-1)'} = \lim_{x \to 1} \frac{\frac{1}{x}}{1} = \lim_{x \to 1} \frac{1}{x} = 1 \text{, έχουμε:} \]
        \[ \lim_{x \to 1} f(x) = \lim_{x \to 1} g(x) \cdot \lim_{x \to 1} \frac{\ln x}{x-1} = 0 \cdot 1 = 0. \]
    \end{rlist}
    \item [β)] Το $1$ είναι ρίζα της εξίσωσης $f(x) = 0$, διότι $f(1) = 0$.

    Επιπλέον, επειδή $f'(x) = \sqrt{x^2+1} > 0$ για κάθε $x \in \mathbb{R}$, η $f$ είναι γνησίως αύξουσα στο $\mathbb{R}$ και ως εκ τούτου είναι «1 – 1». Επομένως η εξίσωση $f(x) = 0$ έχει μία το πολύ ρίζα. Συνδυάζοντας τα παραπάνω προκύπτει ότι η εξίσωση $f(x) = 0$ έχει μία ακριβώς ρίζα, το $1$.

    \item [γ)] Η συνάρτηση $f$ είναι γνησίως αύξουσα στο $\mathbb{R}$ και $f(1) = 0$. Επομένως, έχουμε:
    \begin{itemize}
        \item $x > 1 \Leftrightarrow f(x) > f(1) \Leftrightarrow f(x) > 0$
        \item $x < 1 \Leftrightarrow f(x) < f(1) \Leftrightarrow f(x) < 0$
    \end{itemize}

    \item [δ)] Είναι $f(x) \leq 0$, $x \in [0,1]$. Επομένως
    \begin{align*}
        E &= \int_{0}^{1} |f(x)| dx = \int_{0}^{1} -f(x) dx = -\int_{0}^{1} f(x) dx = -\int_{0}^{1} (x)' f(x) dx \\
        &= - \left([xf(x)]_0^1 - \int_{0}^{1} x f'(x) dx \right) = - \left([xf(x)]_0^1 - \int_{0}^{1} x \sqrt{x^2+1} dx \right) \\
        &= -[xf(x)]_0^1 + \frac{1}{2} \int_{0}^{1} 2x(x^2+1)^{\frac{1}{2}} dx \\
        &= -[xf(x)]_0^1 + \frac{1}{2} \int_{0}^{1} (x^2+1)^{\frac{1}{2}} (x^2+1)' dx \\
        &= -[xf(x)]_0^1 + \frac{1}{2} \left[ \frac{(x^2+1)^{\frac{3}{2}}}{\frac{3}{2}} \right]_0^1 = -(1 \cdot f(1) - 0 \cdot f(0)) + \frac{1}{2} \cdot \frac{2}{3} \left[(x^2+1)^{\frac{3}{2}}\right]_0^1 \\
        &= \frac{1}{3} (2\sqrt{2}-1).
    \end{align*}
    Σχόλιο: Εναλλακτικά για τον υπολογισμό του $\int_{0}^{1} x \sqrt{x^2+1} dx$ μπορούμε να χρησιμοποιήσουμε την μέθοδο αντικατάστασης, θέτοντας $u = \sqrt{x^2+1}$.
\end{alist}